\section{Unfolding and the problem of finding constant series}
\label{sec:unfolding}

\subsection{What a volume computation determines}

The problem is to integrate a sum over curves while the underlying
surface also varies. Directly, this asks us to keep track of all the
curves on every surface in moduli space. \emph{Unfolding} changes the
order of description: we integrate over surfaces equipped with one
chosen curve, and evaluate just the term belonging to that curve.
The sum disappears because the choice of curve has become part of
the integration domain. An elementary example shows exactly how
this trade works.

\paragraph{A toy model: unfolding the circle to the line.}
Let $f:\RR\to[0,\infty)$ be integrable, and form the periodic sum
\[
 P_f(x)=\sum_{n\in\ZZ}f(x+n).
\]
Shifting $x$ by one merely reindexes the sum, so $P_f$ is a
function on the circle $\RR/\ZZ$. A fundamental interval for
that circle is $[0,1)$: each point of the circle has one
representative there. By nonnegativity, Tonelli's theorem allows
us to interchange sum and integral:
\[
 \begin{aligned}
 \int_0^1P_f(x)\,\dd x
 &=\sum_{n\in\ZZ}\int_0^1f(x+n)\,\dd x\\
 &=\sum_{n\in\ZZ}\int_n^{n+1}f(u)\,\dd u
 =\int_{\RR}f(u)\,\dd u.
 \end{aligned}
\]
In the second line we used $u=x+n$. The translated intervals
$[n,n+1)$ tile the real line, so they join into a single
integration domain. This is unfolding: one integral of a
sum on a quotient becomes one integral of a single function on
a larger space. For a signed or complex-valued $f$, the same
argument works if $\int_{\RR}|f|<\infty$.

\begin{figure}[htbp]
\centering
\begin{tikzpicture}[font=\small,>=Latex]
 \draw[->,mpgray] (-4.6,0)--(5,0) node[right] {$u$};
 \foreach \x/\lab in {-4/-2,-2/-1,0/0,2/1,4/2}
   {\draw (\x,-.1)--(\x,.1);\node[below] at (\x,-.1) {$\lab$};}
 \draw[mpblue,line width=2pt] (0,0)--(2,0);
 \draw[mpteal,line width=2pt] (-2,0)--(0,0);
 \draw[mpochre,line width=2pt] (2,0)--(4,0);
 \draw[->,mpteal] (1,.3) to[bend right=35] (-1,.3);
 \draw[->,mpochre] (1,.3) to[bend left=35] (3,.3);
 \node[above,text=mpblue] at (1,1.1) {$x\in[0,1)$};
 \node[above,text=mpteal] at (-1,.75) {$u=x-1$};
 \node[above,text=mpochre] at (3,.75) {$u=x+1$};
 \node[below] at (-4,-.6) {$\cdots$};
 \node[below] at (4,-.6) {$\cdots$};
\end{tikzpicture}
\caption*{Unfolding replaces the copies $f(x+n)$ on one unit interval
by $f(u)$ on all translated intervals $[n,n+1)$. The intervals tile
the line, with each endpoint assigned to the interval on its right.
The picture concerns integration domains; it does not assert that
the periodic sum has a constant value.}
\end{figure}

Take $f(u)=e^{-|u|}$. For $0\le x\le1$, separate $n\ge0$ from
$n\le-1$ and sum the two geometric series:
\[
 P_f(x)=\frac{e^{-x}+e^{x-1}}{1-e^{-1}},
 \qquad \int_0^1P_f(x)\,\dd x=2.
\]
Knowing the integral has not determined the function:
\[
 P_f(0)=\frac{1+e^{-1}}{1-e^{-1}},\qquad
 P_f(\tfrac12)=\frac{2e^{-1/2}}{1-e^{-1}},
\]
and these values differ by
$(1-e^{-1/2})^2/(1-e^{-1})>0$.
Thus unfolding computes the mean, here $2$, without making
the periodic sum constant.

A contrasting example is the tent function
$f(u)=\max(1-|u|,0)$. For $0\le x\le1$ only the terms with
$n=0,-1$ can be nonzero, and their sum is
$(1-x)+x=1$. Its translates form a \emph{partition of unity}:
their values add to one at every point. This extra pointwise
fact explains constancy; integrating it also gives the value
of the integral. McShane's identity will supply a geometric
partition playing a comparable role.

\paragraph{What replaces the circle and the line?}
Fix an oriented surface and its boundary lengths.
Teichm\"uller space $\T$ consists of its marked hyperbolic
structures: the marking identifies curves with curves on a
fixed reference surface. The mapping class group $\Gamma$
consists of orientation-preserving diffeomorphisms modulo
isotopy, with the prescribed boundary conventions. Its action
changes the marking. The quotient $\Gamma\backslash\T$ is
moduli space, where the marking has been forgotten.

Choose an essential, nonperipheral curve class $\gamma$: it cannot
be contracted to a point or moved into a puncture or boundary. Its orbit
$\Gamma\gamma$ consists of all classes obtained by changing
the marking, each counted once. Write $\ell_\delta(X)$ for
the length of the geodesic representative of $\delta$ on $X$.
The subgroup $\Stab(\gamma)$ preserves $\gamma$.
Consequently $\Stab(\gamma)\backslash\T$ remembers a surface
\emph{together with a distinguished curve in this orbit}.
The map to $\Gamma\backslash\T$ forgets that curve.
The different choices over a surface are precisely what the
orbit sum counts, with automorphisms handled by orbifold
integration. Curves of equal length are still distinct
summands when their classes are different.

Use the same Weil--Petersson (WP) orbifold measure $\dV$ on
both quotients. For a measurable nonnegative function $F$,
or one for which
$\int_{\Gamma\backslash\T}\sum_\delta|F(\ell_\delta)|\,\dV<\infty$,
the basic unfolding identity is
\begin{equation}
 \int_{\Gamma\backslash\T}
 \sum_{\delta\in\Gamma\gamma}F(\ell_\delta(X))\,\dV(X)
 =\int_{\Stab(\gamma)\backslash\T}
             F(\ell_\gamma(X))\,\dV(X).
 \label{eq:unfold}
\end{equation}
To prove it, take a fundamental domain for $\Gamma$ and
integrate each summand over that domain. Change the marking
in each term so that the chosen curve becomes $\gamma$.
The resulting translates assemble into a fundamental domain
for $\Stab(\gamma)$, just as the translated unit intervals
assembled into $\RR$. The stabilizer identifies exactly the
repetitions that would otherwise count the same chosen curve
again. Ineffective isotropy is retained on both sides: even symmetries
that act trivially contribute the reciprocal of their group order.
For nonnegative functions either integral may be infinite.

The same argument applies to a specified multicurve, with
$\Stab(\gamma)$ preserving its entire decoration, including
any orientation, ordering or component weights, and with the
same length statistic on both sides. Simplicity is not needed
for this change of integration domain.

\paragraph{Why remembering a simple curve helps.}
If the chosen curve is simple, meaning it has no self-intersections,
cutting along it produces two
geodesic boundary circles of the same length $x$.
To reconstruct the surface we must specify the cut surface
and how far one boundary is shifted before regluing it to the
other. That shift is the Fenchel--Nielsen twist $\tau$,
measured in length units. A full Dehn twist changes $\tau$
to $\tau+x$, so a fundamental twist interval is $0\le\tau<x$.
The WP length--twist factor is $\dd x\,\dd\tau$; integrating
out the twist therefore gives
\[
 \int_0^x\dd\tau\,\dd x=x\,\dd x.
\]
This is the origin of the length factor in the unfolded
integral. It is the size of the gluing parameter, rather than
an extra weight chosen for the curve.

For a disjoint simple multicurve, repeat this construction
along each component. The result has factors
$\prod_j x_j\,\dd x_j$ and the WP volumes of the cut components.
For example, cutting a nonseparating curve on a genus-$g$
surface with $n$ boundaries leaves a genus-$(g-1)$ surface
with $n+2$ boundaries, the two new lengths both being $x$.
A separating cut leaves two surfaces, so their volume factors
multiply. Finite symmetries must still be accounted for:
unoriented cuffs or indistinguishable components must not be
accidentally counted twice. Mirzakhani's integration formula
makes these factors precise~\cite[Theorem 7.1 and Lemma 7.3]{MirzVolumes}.
Its symmetry and one-handle factors must be converted using
the volume convention for each component; they are not extra
factors to append to an already stack-normalized quotient.

\paragraph{An integral is an average; an identity needs more.}
Suppose a series $S(X)$ is independently known to be constant,
say $C$, and $0<V_{g,n}<\infty$ in the chosen normalization.
Then
\begin{equation}
 C=\frac{1}{V_{g,n}}
       \int_{\M_{g,n}}S(X)\,\dV(X),
 \label{eq:recoverconstant}
\end{equation}
and unfolding can calculate its value. Without an independent
constancy argument the same quotient is only the mean of $S$,
as the exponential toy example shows. Conversely, if a pointwise
identity supplies a known nonzero $C$, evaluating its integral
can determine the unknown volume $V_{g,n}$. This second use is
the starting point of Mirzakhani's volume recursion.

\paragraph{A complete surface example: the punctured torus.}
Let $X$ be a complete finite-area hyperbolic torus with one cusp.
All essential unoriented simple closed curves belong to a
single mapping class group orbit. Cutting along a chosen one,
of length $x$, leaves a pair of pants with boundary lengths
$(0,x,x)$, where length zero denotes the cusp. A hyperbolic
pair of pants is determined by its three boundary lengths.
There is therefore no further moduli integral for the cut
surface: the remaining continuous parameters are just $x>0$
and the regluing twist $\tau$.

The effective stabilizer quotient has coordinates
$x>0$ and $0\le\tau<x$. In the stack convention of
Section~\ref{sec:torussphere}, the torus's generic elliptic
involution has order two and acts trivially on Teichm\"uller
space. Orbifold integration retains its weight $1/2$.
Thus the unfolded integral is
$\frac12\int_0^\infty\int_0^x f(x)\,\dd\tau\,\dd x$, giving
\begin{equation}
 \int_{\M_{1,1}}\sum_{\alpha\text{ simple}}f(\ell_\alpha)\,\dV
       =\frac12\int_0^\infty x f(x)\,\dd x.
 \label{eq:torusunfold}
\end{equation}
The sum is over unoriented, nonperipheral simple curve classes.
The factor $1/2$ occurs once; it is not a second quotient by
cuff exchange.

Choose $f(x)=(1+e^x)^{-1}$. The many-curve integral has become
an elementary one-variable integral:
\[
 \int_0^\infty\frac{x\,\dd x}{1+e^x}
 =\sum_{m\ge1}(-1)^{m-1}\int_0^\infty xe^{-mx}\,\dd x
 =\sum_{m\ge1}\frac{(-1)^{m-1}}{m^2}
 =\frac{\pi^2}{12}.
\]
Here $(1+e^x)^{-1}=\sum_{m\ge1}(-1)^{m-1}e^{-mx}$ for $x>0$.
Termwise integration is justified by
$\sum_m\int_0^\infty xe^{-mx}\,\dd x=\sum_m m^{-2}<\infty$.
For a finite numerical check, the integral beyond $R$ is at most
$(R+1)e^{-R}$, and the alternating sum through $m=M$ has error
at most $(M+1)^{-2}$.
Equation~\eqref{eq:torusunfold} therefore gives $\pi^2/24$
for the integral over moduli space.

Why choose this particular $f$? Because McShane's theorem says
that its curve sum is pointwise constant~\cite{McShane}:
\begin{equation}
 \sum_{\alpha\text{ simple, unoriented}}
                 \frac{1}{1+e^{\ell_\alpha(X)}}=\frac12.
 \label{eq:mcshane}
\end{equation}
The geometric idea is a partition at the cusp on a
\emph{single fixed surface}. Take a small horocycle around
the cusp and normalize its length measure to have total mass
one. Its points specify geodesic rays leaving the cusp.
The gap construction groups the complementary intervals by
simple curves; the exceptional set has measure zero.
For a curve of length $x$, the associated gaps have combined
normalized measure $2/(1+e^x)$. This width is computed in the
pair of pants obtained by cutting along that curve.
Adding the gap widths gives one, hence \eqref{eq:mcshane}.
This describes the mechanism, not a proof of the required
gap classification or null-set theorem; see
\cite[Sections 3, 5, and 9]{McShanePants}.
As with the tent function, the partition supplies the
pointwise statement that integration alone cannot supply.

We can now use the two ingredients in either direction.
Integrating McShane's identity and using the computed integral
gives
\[
 \frac12 V_{1,1}=\frac{\pi^2}{24},
 \qquad V_{1,1}=\frac{\pi^2}{12}.
\]
Alternatively, if the volume is already known from the
tetrahedral calculation and torus--sphere correspondence,
unfolding checks that the mean of the McShane sum is $1/2$.
Its constancy still comes from McShane's theorem.
The cusp partition and unfolding do different jobs:
one sums gaps on one surface; the other integrates a curve
sum over all surfaces.

\paragraph{How the example becomes a volume recursion.}
For a surface with positive boundary lengths, Mirzakhani's
generalized McShane identity partitions a distinguished
boundary of length $L_1$ by gap functions associated with
embedded pairs of pants. Its right-hand side is $L_1$.
Integration makes that side $L_1V_{g,n}(L)$.
On the other side, group the pants by topological type and
unfold each orbit sum. Removing the chosen pants leaves
surfaces of smaller complexity. Their volumes, together with
the explicitly known gap kernels and the length--twist
factors, evaluate the unfolded terms.
Consequently the unknown volume is related to lower-complexity
volumes; this is a recursion, not just another integral
representation~\cite[Section 8]{MirzVolumes}.
The punctured-torus example is especially simple because its
cut surface is already a pair of pants, whose fixed-boundary
moduli space is a point.

\subsection{Operations that preserve constancy}

The pointwise mechanism should remain visible when building new
identities. If $S_1(X)=c_1$ and $S_2(X)=c_2$, their product is
$c_1c_2$, but expanding it sums over every pair of indexing curves.
Some pairs intersect. Discarding those terms changes the function
and needs a new argument. This is a useful first check on proposed
multicurve identities.

Finite linear combinations of constant series preserve constancy.
Parameter integrals do too when the constant is integrable and Fubini's
theorem applies to the joint sum and integral, or when all terms are
nonnegative and Tonelli's theorem applies. For differentiation on an
identified parameter family, one sufficient hypothesis is local uniform
absolute convergence of both the series and its first derivatives.
Uniform convergence of the original series alone is insufficient.
Products give constant double series when Tonelli's theorem or absolute
convergence applies. These operations construct kernels together with
their required convergence hypotheses.

Their geometric interpretation needs to be retained. A product usually
indexes pairs of curves that may intersect; it is not automatically a
disjoint-multicurve sum. A parameter derivative also differentiates lengths
if the metric varies with that parameter. In particular, differentiating
a constant series with respect to the metric gives the one-form identity
\begin{equation}
 \sum_\gamma F'(\ell_\gamma)\,\dd\ell_\gamma=0,
 \label{eq:derivativeidentity}
\end{equation}
under the stated local $C^1$ control, for a kernel with no additional
explicit metric dependence. It does not generally give
another scalar identity by differentiating only the displayed weight.

Two possible tests for a proposed constant series follow. One can try to
show that~\eqref{eq:derivativeidentity} vanishes on a spanning set of tangent
directions at every point; a $C^1$ series is then constant on each
connected component. Alternatively, for real-valued $S\in L^2$ on a
finite positive-volume moduli space, let $\overline S$ be its mean. Then
\begin{equation}
 \int_{\M}(S-\overline S)^2\,\dV=0
 \label{eq:variance}
\end{equation}
proves constancy almost everywhere, and everywhere if $S$ is continuous,
since WP measure has full support. For a complex-valued series use
$|S-\overline S|^2$ instead of the square.
The second-moment calculation involves pairs of curves, including their
intersection patterns. It can be substantially harder than the first
unfolding integral.

Thus a useful search for further identities should ask for more than
formal closure: a natural index set, a manageable positive weight, a
geometric partition or telescoping mechanism, and a convergence proof.
The constructions in Sections~\ref{sec:figures} and~\ref{sec:exacttwist}
illustrate these requirements.

\subsection{Other established constant series}

There are also constant series with substantially different kernels.
For example, write
\[
 \mathcal L(x)=\operatorname{Li}_2(x)+\tfrac12\log x\log(1-x)
\]
for the Rogers dilogarithm on $(0,1)$, extended continuously to its
endpoints, and define
\begin{align*}
 D(r)={}&\mathcal L\bigl(\sech^2(r/2)\bigr)\\
 &+2\mathcal L\left(\frac{1+e^{-r}}2\right)
  -2\mathcal L\left(\frac{1-e^{-r}}2\right).
\end{align*}
On a complete finite-area once-punctured torus, Hu--Tan's identity is
\begin{equation}
 \sum_{\alpha\text{ simple, unoriented}}D(\ell_\alpha)=\frac{\pi^2}{2}.
 \label{eq:dilogtorus}
\end{equation}
Here the index set is the nonperipheral simple spectrum with its curve
multiplicities~\cite[Theorem 2, equation (2)]{HuTan}. Their Theorem 1
treats a positive geodesic boundary with a different, boundary-dependent
kernel; the displayed $D$ is the cusped specialization.

For a closed oriented surface of genus $g\ge2$, Luo--Tan give explicit
dilogarithmic weights $f(P),g(T)$ on properly embedded geometric pairs of
pants and one-holed tori, each subsurface counted once,
such that
\begin{equation}
 \sum_P f(P)+\sum_T g(T)=8\pi^2(g-1)=\operatorname{Vol}(T^1X).
 \label{eq:luotan}
\end{equation}
Their construction partitions unit tangent directions up to a null set
\cite[Theorem 1.1 and Section 1.2]{LuoTan}. Angular measure $2\pi$ on
each unit tangent circle, together with area $4\pi(g-1)$, explains the
normalization on the right.
It illustrates another source of constant series: the invariant measure
being partitioned need not live only on a horocycle or a geodesic boundary.
The precise, longer subsurface weights are given in that reference.

\paragraph{Computational support.}
Group G1 of \path{experiments/quantum_geodesics/README.md} checks
the finite torus normalization and Fermi-integral formulas, with the
specified finite cutoffs and analytic tail. Its numerical comparisons
do not prove a quotient-volume identity or constancy of a general
surface sum. The result path is
\path{artifacts/quantum_geodesics/result.json};
\path{book/claims/sections17-22.json} records the corresponding claims
and the separate scope of the analytic and computational evidence.
